%===========================================================================
% MISSED CFA LEVEL II CONCEPTS - NEW SLIDES ONLY
% These concepts were missing from the original presentation
%===========================================================================

\documentclass[aspectratio=169,10pt]{beamer}
\usetheme{Madrid}
\usecolortheme{default}
\usepackage{amsmath, amssymb, amsfonts}
\usepackage{booktabs}
\usepackage{array}
\usepackage{multirow}
\usepackage{bm}
\usepackage{xcolor}
\usepackage{tikz}
\usepackage{enumitem}
\usepackage{caption}
\usepackage{subcaption}
\usepackage{adjustbox}
\usepackage{multicol}
\usetikzlibrary{shapes,arrows,positioning,calc}

% Define colors
\definecolor{navyblue}{RGB}{0,53,102}
\definecolor{crimson}{RGB}{165,42,42}
\definecolor{darkgreen}{RGB}{0,100,0}
\definecolor{gold}{RGB}{184,134,11}
\definecolor{lightgray}{RGB}{240,240,240}
\definecolor{darkgray}{RGB}{100,100,100}
\definecolor{cfablue}{RGB}{0,70,150}
\definecolor{answerbox}{RGB}{245,245,245}

% Title formatting
\setbeamercolor{title}{fg=navyblue}
\setbeamercolor{frametitle}{fg=navyblue}
\setbeamercolor{structure}{fg=navyblue}
\setbeamercolor{block title}{bg=navyblue!20, fg=navyblue}
\setbeamercolor{alertblock title}{bg=crimson!20, fg=crimson}
\setbeamercolor{exampleblock title}{bg=darkgreen!20, fg=darkgreen}

% Smaller fonts
\renewcommand{\tiny}{\fontsize{4pt}{5pt}\selectfont}
\renewcommand{\scriptsize}{\fontsize{5pt}{6pt}\selectfont}
\renewcommand{\footnotesize}{\fontsize{6pt}{7pt}\selectfont}
\renewcommand{\small}{\fontsize{7pt}{8.5pt}\selectfont}
\renewcommand{\normalsize}{\fontsize{8pt}{10pt}\selectfont}
\renewcommand{\large}{\fontsize{9pt}{11pt}\selectfont}
\renewcommand{\Large}{\fontsize{10pt}{12pt}\selectfont}
\renewcommand{\LARGE}{\fontsize{11pt}{13pt}\selectfont}
\renewcommand{\huge}{\fontsize{12pt}{14pt}\selectfont}
\renewcommand{\Huge}{\fontsize{14pt}{16pt}\selectfont}

% Custom commands
\newcommand{\question}[1]{{\large\bfseries\color{crimson}#1}}
\newcommand{\subquestion}[1]{{\bfseries\color{darkgreen}#1}}
\newcommand{\concept}[1]{\textcolor{navyblue}{\textbf{#1}}}
\newcommand{\highlight}[1]{\textcolor{crimson}{\textbf{#1}}}
\newcommand{\boxedanswer}[1]{\fcolorbox{crimson}{answerbox}{\parbox{0.9\linewidth}{\small #1}}}

\begin{document}
	
	%===========================================================================
	% MISSED CONCEPT 1: MODERN PORTFOLIO THEORY (MPT) TO BEHAVIORAL FINANCE
	%===========================================================================
	
	\begin{frame}{Classical to Contemporary: MPT $\rightarrow$ Behavioral Finance}
		\question{Classical Foundation: Modern Portfolio Theory (Markowitz, 1952)}
		\small
		\begin{itemize}
			\item Risk should not be viewed in isolation, but as part of a diversified portfolio
			\item Assumes \textbf{Homo economicus}: perfectly rational actor
			\item Investors maximize expected utility while avoiding unnecessary risk
			\item Markets are highly efficient
			\item \textbf{Key Formula:} Portfolio variance = $\sum w_i^2\sigma_i^2 + 2\sum\sum w_iw_j\sigma_i\sigma_j\rho_{ij}$
		\end{itemize}
		
		\vspace{0.1cm}
		\question{Contemporary Successor: Behavioral Finance \& Neuroeconomics}
		\small
		\begin{itemize}
			\item Real-world crises revealed humans deviate from strict rationality
			\item \textbf{Neuroeconomics}: Financial decision-making fluctuates across lifespan (Kuhnen, 2025)
			\item Risk tolerance is inherently volatile and subjective
			\item Cognitive \& emotional factors influence decisions
			\item \textbf{Key Insight:} Risk aversion changes with age, experience, and context
		\end{itemize}
	\end{frame}
	
	%===========================================================================
	% MISSED CONCEPT 2: MPT BUSINESS APPLICATIONS
	%===========================================================================
	
	\begin{frame}{MPT to Behavioral Finance: Business Applications}
		\question{Real-World Applications}
		\small
		
		\textbf{1. Wealth Management \& Product Design}
		\begin{itemize}
			\item Vanguard, Fidelity use MPT for baseline asset allocations
			\item Apply neuroeconomic insights to client relationship management
			\item Target-date funds designed for aging populations
			\item Automated "nudge" algorithms prevent panic-selling
			\item Programmable ledgers and tokenization build automated guardrails (BIS, 2025)
		\end{itemize}
		
		\vspace{0.1cm}
		\textbf{2. Enterprise Risk Management (ERM)}
		\begin{itemize}
			\item Executives prone to heuristic biases (overconfidence, loss aversion)
			\item Destructive M\&A overbidding, poor capital structuring
			\item ERM forces "portfolio view of risk" (Beasley \& AICPA, 2025)
			\item Oversight committees as structural counterweight to cognitive bias
			\item Corporate risk-taking aligned with rational ideals of MPT
		\end{itemize}
	\end{frame}
	
	%===========================================================================
	% MISSED CONCEPT 3: EMH TO ADAPTIVE MARKETS HYPOTHESIS (AMH)
	%===========================================================================
	
	\begin{frame}{Classical to Contemporary: EMH $\rightarrow$ AMH}
		\question{Classical: Efficient Market Hypothesis (Fama, 1970)}
		\small
		\begin{itemize}
			\item Stock prices reflect all available information
			\item Impossible to consistently achieve above-market returns
			\item Investors are rational
			\item Led to growth of \textbf{passive investing} (Vanguard, BlackRock)
		\end{itemize}
		
		\vspace{0.1cm}
		\question{Contemporary: Adaptive Markets Hypothesis (Lo, 2004)}
		\small
		\begin{itemize}
			\item Markets function like ecosystems - investors learn, compete, adapt
			\item Efficiency fluctuates over time
			\item Market inefficiency follows \textbf{predictable patterns}
			\item Behavioral factors are permanent features of markets
			\item \textbf{Key Insight:} "Failures follow predictable patterns" - Barberis (2014)
		\end{itemize}
		
		\vspace{0.1cm}
		\question{Key Difference:}
		\small
		\begin{itemize}
			\item EMH: Information is always fully reflected in prices
			\item AMH: Information processing depends on market conditions and investor psychology
		\end{itemize}
	\end{frame}
	
	%===========================================================================
	% MISSED CONCEPT 4: AMH REAL-WORLD EVIDENCE
	%===========================================================================
	
	\begin{frame}{AMH: Real-World Evidence}
		\question{Example 1: 2008 Global Financial Crisis}
		\small
		\begin{itemize}
			\item Investors viewed mortgage-backed securities as safe
			\item Prices should have reflected all information (EMH view)
			\item Widespread mispricing of risk demonstrated slow adaptation
			\item AMH: Investors learn through experience, respond imperfectly
			\item Adaptation occurred after significant market disruption
		\end{itemize}
		
		\vspace{0.1cm}
		\question{Example 2: Algorithmic Trading}
		\small
		\begin{itemize}
			\item Renaissance Technologies, Two Sigma continuously modify models
			\item Success depends on ability to evolve
			\item Reflects adaptive principles of AMH
			\item Dynamic environments require continuous adjustment and innovation
		\end{itemize}
		
		\vspace{0.1cm}
		\question{Example 3: Peloton Interactive (2019-2023)}
		\small
		\begin{itemize}
			\item IPO at \$29, peaked at \$167 (market cap \$49B)
			\item Never posted operating profit
			\item Cognitive biases: herding, extrapolation, loss aversion
			\item By late 2022: trading in single digits, \$2.8B loss
			\item EMH failed; AMH explained predictable patterns
		\end{itemize}
	\end{frame}
	
	%===========================================================================
	% MISSED CONCEPT 5: CAPM TO MULTIFACTOR MODELS
	%===========================================================================
	
	\begin{frame}{Classical to Contemporary: CAPM $\rightarrow$ Multifactor}
		\question{Classical: Capital Asset Pricing Model (Sharpe, 1964; Lintner, 1965)}
		\small
		\[
		E(R_i) = R_f + \beta_i[E(R_m) - R_f]
		\]
		\begin{itemize}
			\item Links expected return to \textbf{single beta} (market risk)
			\item Assumes one factor explains all returns
			\item Foundation for WACC and cost of equity calculations
			\item \textbf{Critique:} Misses size, value, profitability, investment effects
		\end{itemize}
		
		\vspace{0.1cm}
		\question{Contemporary: Fama-French Multifactor Models}
		\small
		
		\textbf{3-Factor Model (Fama \& French, 1993):}
		\[
		E(R_i) = R_f + \beta_m(R_m - R_f) + \beta_{SMB}SMB + \beta_{HML}HML
		\]
		
		\textbf{5-Factor Model (Fama \& French, 2015):}
		\[
		E(R_i) = R_f + \beta_m(R_m - R_f) + \beta_{SMB}SMB + \beta_{HML}HML + \beta_{RMW}RMW + \beta_{CMA}CMA
		\]
		\begin{itemize}
			\item SMB = Size (Small minus Big)
			\item HML = Value (High minus Low)
			\item RMW = Profitability (Robust minus Weak)
			\item CMA = Investment (Conservative minus Aggressive)
		\end{itemize}
	\end{frame}
	
	%===========================================================================
	% MISSED CONCEPT 6: MULTIFACTOR BUSINESS APPLICATIONS
	%===========================================================================
	
	\begin{frame}{Multifactor Models: Business Applications}
		\question{Real-World Applications}
		\small
		
		\textbf{1. Fund Management}
		\begin{itemize}
			\item Fund companies offer factor-tilted index funds
			\item Based on multifactor ideas (size, value, profitability, investment)
			\item Pensions evaluate managers against several betas
			\item Not just market beta anymore
		\end{itemize}
		
		\vspace{0.1cm}
		\textbf{2. Corporate Finance}
		\begin{itemize}
			\item Finance teams start with CAPM for WACC
			\item Compare with multifactor expected returns
			\item Set hurdle rates using multiple risk factors
			\item More accurate cost of capital estimates
		\end{itemize}
		
		\vspace{0.1cm}
		\textbf{3. Factor Investing}
		\begin{itemize}
			\item Smart beta strategies
			\item Factor tilts for enhanced returns
			\item Risk factor diversification
			\item Evidence: Size, value, profitability premiums persist
		\end{itemize}
	\end{frame}
	
	%===========================================================================
	% MISSED CONCEPT 7: BEHAVIORAL FINANCE - CORE PRINCIPLES
	%===========================================================================
	
	\begin{frame}{Behavioral Finance: Core Principles}
		\question{Core Principles (Barberis, 2014)}
		\small
		
		\textbf{1. Cognitive Biases}
		\begin{itemize}
			\item Overconfidence: Overestimate ability and knowledge
			\item Anchoring: Fixate on initial information
			\item Confirmation Bias: Seek confirming evidence
			\item Availability Bias: Overweight recent events
			\item Loss Aversion: Fear losses more than value gains
		\end{itemize}
		
		\vspace{0.1cm}
		\textbf{2. Prospect Theory (Kahneman \& Tversky, 1979)}
		\begin{itemize}
			\item People value gains and losses differently
			\item Losses hurt about twice as much as gains feel good
			\item Decisions based on reference points, not absolute wealth
			\item Probability weighting: over-weight small probabilities, under-weight large
		\end{itemize}
		
		\vspace{0.1cm}
		\textbf{3. Key Insight}
		\begin{itemize}
			\item Market inefficiency is not random
			\item Failures follow predictable patterns
			\item Behavioral Finance explains what EMH cannot
		\end{itemize}
	\end{frame}
	
	%===========================================================================
	% MISSED CONCEPT 8: NEUROECONOMICS
	%===========================================================================
	
	\begin{frame}{Neuroeconomics: The Biological Foundation}
		\question{Neuroeconomics (Kuhnen, 2025)}
		\small
		
		\textbf{Key Findings:}
		\begin{itemize}
			\item Financial decision-making varies across lifespan
			\item Neurological changes affect risk processing
			\item Emotional states influence financial choices
			\item Cognitive aging impacts investment decisions
			\item Same cognitive machinery operates for individuals and fund managers
		\end{itemize}
		
		\vspace{0.1cm}
		\textbf{Implications for Business:}
		\begin{itemize}
			\item Institutional title does not grant immunity to biases
			\item Risk tolerance is inherently volatile and subjective
			\item Organizations need structural counterweights
			\item ERM frameworks, oversight committees, decision protocols
		\end{itemize}
		
		\vspace{0.1cm}
		\textbf{Business Applications:}
		\begin{itemize}
			\item Target-date funds for aging populations
			\item Automated guardrails to prevent panic-selling
			\item Programmable ledgers (BIS, 2025)
			\item Risk committee structures
		\end{itemize}
	\end{frame}
	
	%===========================================================================
	% MISSED CONCEPT 9: BIS NEXT-GENERATION FINANCIAL SYSTEMS
	%===========================================================================
	
	\begin{frame}{BIS 2025: Next-Generation Financial Systems}
		\question{Bank for International Settlements (2025)}
		\small
		
		\textbf{1. Next-Generation Monetary Systems}
		\begin{itemize}
			\item Central Bank Digital Currencies (CBDCs) emerging
			\item Create new risks and opportunities
			\item Risk management must adapt to new technologies
		\end{itemize}
		
		\vspace{0.1cm}
		\textbf{2. Tokenization and Programmable Ledgers}
		\begin{itemize}
			\item Smart contracts embed risk management rules
			\item Programmable money enables automated compliance
			\item Reduces operational risk
			\item Automated guardrails mitigate human behavioral risks
		\end{itemize}
		
		\vspace{0.1cm}
		\textbf{3. Implications for Financial Firms}
		\begin{itemize}
			\item Invest in technology infrastructure
			\item Develop AI-enabled risk management
			\item Ensure compliance with evolving regulations
			\item Maintain human judgment in decision-making
		\end{itemize}
	\end{frame}
	
	%===========================================================================
	% MISSED CONCEPT 10: ERM AS STRUCTURAL COUNTERWEIGHT
	%===========================================================================
	
	\begin{frame}{ERM as Structural Counterweight to Biases}
		\question{Enterprise Risk Management (Beasley \& AICPA, 2025)}
		\small
		
		\textbf{Why ERM is Needed:}
		\begin{itemize}
			\item Executives prone to heuristic biases (overconfidence, loss aversion)
			\item Cognitive biases lead to destructive decisions
			\item Individual judgment subject to same biases as retail investors
		\end{itemize}
		
		\vspace{0.1cm}
		\textbf{How ERM Counteracts Biases:}
		\begin{itemize}
			\item \textbf{Independent Oversight:} Separate risk function, board access
			\item \textbf{Scenario Planning:} Test tail risk scenarios, challenge assumptions
			\item \textbf{Diverse Perspectives:} External experts, dissenting views
			\item \textbf{Decision Protocols:} Structured approvals, documentation
			\item \textbf{Post-Mortem Analysis:} Review, learn, improve
			\item \textbf{Portfolio View of Risk:} Forces holistic risk assessment
		\end{itemize}
		
		\vspace{0.1cm}
		\textbf{Key Insight:}
		\begin{itemize}
			\item Governance structures serve as institutional counterweight
			\item Organizations adopt portfolio views of risk
			\item Risk committee structures, investment policy statements
			\item Assumption stress-testing before major capital decisions
		\end{itemize}
	\end{frame}
	
	%===========================================================================
	% MISSED CONCEPT 11: THEORY EVOLUTION FRAMEWORK
	%===========================================================================
	
	\begin{frame}{Financial Theory Evolution: Summary Framework}
		\question{Classical to Contemporary Evolution}
		\small
		
		\begin{table}[h]
			\centering
			\begin{tabular}{|c|c|c|}
				\hline
				\textbf{Classical Theory} & \textbf{Contemporary Theory} & \textbf{Key Innovation} \\
				\hline
				MPT (Markowitz, 1952) & Behavioral Finance & Biological realism, cognitive biases \\
				EMH (Fama, 1970) & AMH (Lo, 2004) & Adaptation, predictable patterns \\
				CAPM (Sharpe, 1964) & Fama-French (1993, 2015) & Size, value, profitability, investment \\
				Rational Actor & Neuroeconomics & Neurological, emotional factors \\
				\hline
			\end{tabular}
		\end{table}
		
		\vspace{0.1cm}
		\question{Key Themes}
		\small
		\begin{itemize}
			\item Shift from normative to descriptive models
			\item From Homo economicus to realistic human behavior
			\item From static efficiency to adaptive markets
			\item From single factors to multifactor models
			\item From individual rationality to organizational counterweights
		\end{itemize}
		
		\vspace{0.1cm}
		\question{Modern Synthesis}
		\small
		\begin{itemize}
			\item Classical theories provide mathematical architecture
			\item Contemporary theories add behavioral realism
			\item Successful strategies combine both approaches
			\item ERM, governance, and technology mitigate human biases
		\end{itemize}
	\end{frame}
	
	%===========================================================================
	% MISSED CONCEPT 12: PELOTON CASE STUDY - DETAILED
	%===========================================================================
	
	\begin{frame}{Case Study: Peloton Interactive (2019-2023)}
		\question{Behavioral Finance in Action}
		\small
		
		\textbf{Timeline:}
		\begin{itemize}
			\item \textbf{2019 IPO:} \$29/share
			\item \textbf{Early 2021:} Peaked at \$167/share (Market cap \$49B)
			\item \textbf{Late 2022:} Single digits, \$2.8B net loss
		\end{itemize}
		
		\vspace{0.1cm}
		\textbf{Behavioral Biases at Play:}
		\begin{itemize}
			\item \textbf{Herding:} Investors followed the crowd
			\item \textbf{Extrapolation:} Assumed pandemic demand was permanent
			\item \textbf{Loss Aversion:} Panic selling during downturn
			\item \textbf{Overconfidence:} Ignored warning signs
			\item \textbf{Confirmation Bias:} Focused on positive news
		\end{itemize}
		
		\vspace{0.1cm}
		\textbf{Key Insight (Barberis, 2014):}
		\begin{itemize}
			\item Market inefficiency is not random
			\item Failures follow predictable patterns
			\item Cognitive filters produce systematic errors on available data
			\item EMH failed; AMH and Behavioral Finance explained it
		\end{itemize}
	\end{frame}
	
	%===========================================================================
	% MISSED CONCEPT 13: PRACTICAL FINANCE IMPLICATIONS
	%===========================================================================
	
	\begin{frame}{Practical Implications for Financial Managers}
		\question{What This Means for Practice}
		\small
		
		\textbf{1. Capital Allocation}
		\begin{itemize}
			\item Account for predictable behavioral responses
			\item Don't assume markets always process information efficiently
			\item Use multifactor models for better cost of capital estimates
		\end{itemize}
		
		\vspace{0.1cm}
		\textbf{2. Investor Communication}
		\begin{itemize}
			\item Understand cognitive biases of investors
			\item Frame information to minimize behavioral errors
			\item Use multiple channels and repetition
		\end{itemize}
		
		\vspace{0.1cm}
		\textbf{3. Risk Management}
		\begin{itemize}
			\item Implement ERM frameworks
			\item Create structural counterweights to biases
			\item Use automated guardrails (BIS, 2025)
			\item Scenario planning and stress testing
		\end{itemize}
		
		\vspace{0.1cm}
		\textbf{4. Corporate Governance}
		\begin{itemize}
			\item Independent oversight committees
			\item Diverse perspectives on risk decisions
			\item Documentation and structured approval processes
			\item Post-mortem analysis and continuous improvement
		\end{itemize}
	\end{frame}
	
	%===========================================================================
	% MISSED CONCEPT 14: REFERENCES FOR MISSED CONCEPTS
	%===========================================================================
	
	\begin{frame}{References: Missed Concepts}
		\tiny
		\begin{thebibliography}{99}
			\bibitem{barberis2014} Barberis, N. C. (2014). Behavioral finance: Core principles and practical applications. CFA Institute.
			
			\bibitem{beasley2025} Beasley, M. S., \& AICPA. (2025). The state of risk oversight (16th ed.). NC State University.
			
			\bibitem{bis2025} BIS. (2025). Annual economic report 2025: Chapter III – Next-generation monetary and financial system.
			
			\bibitem{fama1970} Fama, E. F. (1970). Efficient capital markets. Journal of Finance, 25(2), 383-417.
			
			\bibitem{fama1993} Fama, E. F., \& French, K. R. (1993). Common risk factors. JFE, 33(1), 3-56.
			
			\bibitem{fama2015} Fama, E. F., \& French, K. R. (2015). A five-factor asset pricing model. JFE, 116(1), 1-22.
			
			\bibitem{kahneman1979} Kahneman, D., \& Tversky, A. (1979). Prospect theory. Econometrica, 47(2), 263-291.
			
			\bibitem{kuhnen2025} Kuhnen, C. M. (2025). Financial decision-making across the lifespan. Annual Review of Financial Economics, 17, 395-410.
			
			\bibitem{lintner1965} Lintner, J. (1965). The valuation of risk assets. Review of Economics and Statistics, 47(1), 13-37.
			
			\bibitem{lo2004} Lo, A. W. (2004). The adaptive markets hypothesis. Journal of Portfolio Management, 30(5), 15-29.
			
			\bibitem{markowitz1952} Markowitz, H. (1952). Portfolio selection. Journal of Finance, 7(1), 77-91.
			
			\bibitem{openstax2022} OpenStax. (2022). Principles of finance. Rice University.
			
			\bibitem{sharpe1964} Sharpe, W. F. (1964). Capital asset prices. Journal of Finance, 19(3), 425-442.
		\end{thebibliography}
	\end{frame}
	
	%===========================================================================
	% END - MISSED CONCEPTS ONLY
	%===========================================================================
	
\end{document}